\documentclass[11pt, aspectratio=169]{beamer}
% \documentclass[11pt,handout]{beamer}
\usepackage[T1]{fontenc}
\usepackage[utf8]{inputenc}
\usepackage{textcomp}
\usepackage{float, afterpage, rotating, graphicx}
\usepackage{epstopdf}
\usepackage{longtable, booktabs, tabularx}
\usepackage{fancyvrb, moreverb, relsize}
\usepackage{eurosym, calc}
\usepackage{amsmath, amssymb, amsfonts, amsthm, bm}
\DeclareMathOperator*{\argmin}{arg\,min}
\usepackage[
    natbib=true,
    bibencoding=inputenc,
    bibstyle=authoryear-ibid,
    citestyle=authoryear-comp,
    maxcitenames=3,
    maxbibnames=10,
    useprefix=false,
    sortcites=true,
    backend=biber
]{biblatex}
\AtBeginDocument{\toggletrue{blx@useprefix}}
\AtBeginBibliography{\togglefalse{blx@useprefix}}
\setlength{\bibitemsep}{1.5ex}
\addbibresource{refs.bib}

\hypersetup{colorlinks=true, linkcolor=black, anchorcolor=black, citecolor=black, filecolor=black, menucolor=black, runcolor=black, urlcolor=black}

\setbeamertemplate{footline}[frame number]
\setbeamertemplate{navigation symbols}{}
\setbeamertemplate{frametitle}{\centering\vspace{1ex}\insertframetitle\par}


\begin{document}

\title{Fréchet Regression and Functional Data}

\author[Mirkan Öcal]
{
{\bf Mirkan Öcal}\\
{\small University of Bonn}\\[1ex]
}


\begin{frame}
    \titlepage
    \note{~}
\end{frame}


\begin{frame}[allowframebreaks]
    \frametitle{Introduction}
    \begin{itemize}
        \item Thesis compares two methods for regression with density responses
        \item First method from the fda literature: functional regression after the
        log quantile density (LQD) transformation
        \parencites{PetersenMüller2016}{KokoszkaEtAl2019}
        \item Second method: global Fréchet regression from the object-oriented data
        analysis literature \parencite{PetersenMüller2019}
    \end{itemize}
    \framebreak
    \begin{itemize}
        \item The space of densities is not a linear space, so directly applying fda
        methods to densities is inappropriate
        \item Results are only guaranteed to lie in the ambient linear space
        ($\mathcal{L}^2$), not in the density space itself
        \item Two ways out: transform to an unconstrained space first
        \parencite{PetersenMüller2016}, or respect the geometry of the density space
        directly \parencite{PetersenMüller2019}
    \end{itemize}
\end{frame}

\begin{frame}[allowframebreaks]
    \frametitle{Transformation Approach}
    \begin{itemize}
        \item Transform each density $f$ to its log quantile density
        $\psi(f)(u) = -\log f(Q(u))$, an (essentially) unconstrained element of
        $\mathcal{L}^2([0,1])$
        \item Perform standard functional linear regression on the transformed
        functions, then map the fitted values back to density space
        \item Densities with differing supports require storing an anchor point
        $Q(0)$ per density and estimating it for predictions
        \parencite{KokoszkaEtAl2019} --- a key practical drawback
    \end{itemize}
    \framebreak
    \begin{itemize}
        \item Fréchet regression instead generalizes the conditional mean to metric
        spaces: $m_\oplus(x) = \argmin_{\omega} \mathbb{E}[d^2(Y, \omega) \mid X = x]$
        \item With the Wasserstein-2 distance, estimation reduces to weighted
        averaging of quantile functions plus a monotonicity projection
        (quadratic programming)
        \item No transformation, no support estimation, results are valid
        distributions by construction
    \end{itemize}
\end{frame}



\begin{frame}[plain]
    \begin{figure}
        \centering
        % \input{../bld/figures/frechet/some_densities.pgf}
        \resizebox{0.9\textwidth}{0.9\textheight}{\input{../bld/figures/frechet/some_densities.pgf}}
        \caption[Several sampled pdfs]{Plot of several sample pdfs.}
        \label{fig:some_densities}
    \end{figure}
\end{frame}

\begin{frame}
    \frametitle{Simulation Results}
    \begin{itemize}
        \item DGP: normal distributions, location and scale linear in the predictor
        (as in \textcite{PetersenMüller2019}); scenarios with directly observed and
        with estimated densities
    \end{itemize}
    \begin{table}
        \centering
        \begin{tabular}{lccc}
            \hline
            Mean ISE & $n = 50$ & $n = 100$ & $n = 200$ \\
            \hline
            Fréchet             & 0.18 & 0.08 & 0.04 \\
            Fréchet, D.-Est.    & 32.31 & 27.98 & 27.38 \\
            LQD                 & 34.49 & 31.19 & 31.88 \\
            LQD, D.-Est.        & 47.08 & 43.39 & 43.43 \\
            \hline
        \end{tabular}
    \end{table}
    \begin{itemize}
        \item Fréchet regression clearly dominates; most of the LQD error stems from
        having to estimate the support location, which the transformation discards
    \end{itemize}
\end{frame}

\begin{frame}
    \frametitle{Additional Scenarios}
    \begin{itemize}
        \item Three further scenarios vary which method is correctly specified
        (200 Monte Carlo repetitions, directly observed densities)
    \end{itemize}
    \begin{table}
        \centering
        \small
        \begin{tabular}{llccc}
            \hline
            Scenario & Method & $n = 50$ & $n = 100$ & $n = 200$ \\
            \hline
            Common support ($\times 10^{-3}$)  & Fréchet & 0.026 & 0.012 & 0.005 \\
                                               & LQD     & 0.360 & 0.327 & 0.304 \\
            Linear in LQD space ($\times 10^{-3}$) & Fréchet & 4.315 & 4.019 & 3.851 \\
                                               & LQD     & 0.330 & 0.210 & 0.084 \\
            Nonlinear                          & Fréchet & 7.485 & 7.205 & 7.099 \\
                                               & LQD     & 28.708 & 30.793 & 30.972 \\
            \hline
        \end{tabular}
    \end{table}
    \begin{itemize}
        \item Common support makes LQD competitive $\Rightarrow$ support estimation
        drives its error; under misspecification, Fréchet degrades gracefully
    \end{itemize}
\end{frame}

\begin{frame}
    \frametitle{Conclusion}
    \begin{itemize}
        \item Fréchet regression adapts to the geometry of the density space and
        avoids transformation and back-transformation steps entirely
        \item The LQD method is competitive only when the information destroyed by
        the transformation (support location) does not matter
        \item The density estimation step dominates the error of both methods when
        densities are not directly observed --- bandwidth choice matters
    \end{itemize}
\end{frame}


\begin{frame}[allowframebreaks]
    \frametitle{References}
    \renewcommand{\bibfont}{\normalfont\footnotesize}
    \printbibliography
\end{frame}

\end{document}
